\chapter*{Executive Summary}
\addcontentsline{toc}{chapter}{Executive Summary}
\markboth{Executive Summary}{Executive Summary}

\textbf{BuildScan~3D} is a mobile-to-cloud system that allows users to turn a short video of an indoor space captured with a smartphone into an accurate 3D model that can be used to measure.

\paragraph{Problem.} The 3D scanning of interiors is usually performed using costly methods that involve either using LiDAR devices which cost more than \$10,000 and require multiple preconditions or using photogrammetry, which requires advanced experience in the sphere as well as much time to analyze the data. Many construction companies, facility managers, and students cannot afford to buy the necessary equipment and utilize modern technology to create highly accurate models.

\paragraph{Research.} We analyzed the existing developments related to the application of Gaussian Splatting for 3D scanning~\cite{kerbl3dgs}, PocketGS~\cite{pocketgs}, and visual-inertial SLAM (VIGS-SLAM~\cite{vigsslam}) and found a significant gap. Thus, no system has combined the use of smartphones equipped with ARCore and cloud computing, allowing receiving accurate 3D models in a fully automatic manner.

\paragraph{Solution.} (1) An Android app records the video with ARCore poses, IMU data and live image-quality and
coverage guidance, and uploads it; (2) Structure-from-Motion with lens self-calibration gives the camera poses,
scaled and levelled with ARCore; where walls have no texture, a sensor-assisted chain recovers the path from the
gyroscope, the accelerometer and point tracks; (3) 3D Gaussian Splatting is trained with depth, normal and planar
priors, a motion-blur model and flat Gaussians; (4) room layout, semantic labels (walls, floor, ceiling, doors,
windows, switches, sockets, lights), a mesh and a floor plan are derived from the model; (5) a web viewer shows the
model with label layers and measurement. Since November one command runs all of it from the video, choosing the
recipe from the recording itself.
% DRAFT 2026-11-28 (Tayyab's agent): Solution / status / risks / plan rewritten from the records. TODO(team): read.

\paragraph{Status at the end of FYP-1.} The whole chain works from the command line on the development PC
(RTX~3070): on the main test room the automatic pipeline gives 25.33\,dB held-out PSNR (aligned) and a room length
1.86\,\% short of the tape, in 106 minutes with no hand-set option (R-136); the web files are 92\,\% smaller than the
trained model at a render difference of 36\,dB (R-139). On a featureless recording the sensor chain reached 24.25\,dB
by hand (R-133); the automatic version is 1.4\,dB lower (R-135). The upload server still runs the older recipe.

\paragraph{Risks.} Processing time (72 to 106 minutes on an RTX~3070 against the 15-minute target on an RTX~4090);
semantics scored only on the room they were tuned on; room layout on recordings that never show the floor; testing on
a handful of rooms recorded by the team.

\paragraph{FYP-2 plan.} The server and the phone run the automatic pipeline end to end; processing time; semantics
evaluated on new rooms; the sensor chain made as good automatically as by hand; user acceptance test with the
construction stakeholder.
% TODO(team): align with the Remaining FYP Plan chapter.
